\subsection{exp21 - ICache 集成}

\subsubsection{Request Buffer 更新条件错误导致缺失请求地址更改}

最初 Request Buffer 在每次状态转换到 LOOKUP 时都更新，导致 Cache Miss 期间 PC 变化会覆盖掉当前正在处理的缺失请求地址。

\noindent
\textbf{现象：}Cache Miss 处理期间，如果 PC 继续变化（如取下一条指令），Request Buffer 中的地址会被更新为新的 PC，导致缺失处理使用错误的地址向总线发起请求，返回的 Cache 行填充到错误的位置。

\noindent
\textbf{分析：}Cache Miss 后，状态机进入 MISS → REPLACE → REFILL 流程，此时应保持 Request Buffer 不变，使用缺失时的地址进行替换和重填。但若 Request Buffer 在任何 LOOKUP 转换时都更新，连续两次命中后的第三次缺失会导致 Buffer 被第二次的地址覆盖。

\begin{lstlisting}[language=Verilog, caption=错误的 Buffer 更新条件]
// 错误：只要进入 LOOKUP 就更新
always @(posedge clk) begin
    if (!resetn) begin
        buf_tag <= {WIDTH_TAG{1'b0}};
        buf_index <= {WIDTH_INDEX{1'b0}};
        // ...
    end else begin
        if (main_next_state == MAIN_LOOKUP) begin
            buf_tag <= tag;
            buf_index <= index;
            // ...
        end
    end
end
\end{lstlisting}

\noindent
\textbf{解决：}只在握手成功（\verb|valid && addr_ok|）时更新 Request Buffer，确保 Miss 处理期间地址不变：

\begin{lstlisting}[language=Verilog, caption=正确的 Buffer 更新条件]
always @(posedge clk) begin
    if (!resetn) begin
        buf_isstore <= 1'b0;
        buf_tag <= {WIDTH_TAG{1'b0}};
        buf_index <= {WIDTH_INDEX{1'b0}};
        buf_offset <= {WIDTH_OFFSET{1'b0}};
        buf_wdata <= 32'h0;
        buf_wstrb <= 4'h0;
    end else if (valid && addr_ok && !write_conflict) begin
        // 只在握手时锁存请求，确保 miss/refill 期间 PC 变化不会污染地址
        buf_isstore <= op;
        buf_tag <= tag;
        buf_index <= index;
        buf_offset <= offset;
        buf_wdata <= wdata;
        buf_wstrb <= wstrb;
    end
end
\end{lstlisting}

\subsubsection{addr\_ok 生成条件不完整}

最初的 \verb|addr_ok| 生成逻辑过于简化，未考虑 \verb|valid| 信号和命中条件。

\begin{lstlisting}[language=Verilog, caption=错误的 addr_ok 生成]
// 错误：未考虑 valid 和 hit 条件
assign addr_ok = (main_state == MAIN_IDLE)
               | (main_state == MAIN_LOOKUP && main_next_state == MAIN_LOOKUP);
\end{lstlisting}

\noindent
\textbf{现象：}流水线控制异常，IDLE 状态下即使没有请求也会置 \verb|addr_ok|，导致流水线错误推进；LOOKUP 状态下即使miss也会在转换前短暂置 \verb|addr_ok|。

\noindent
\textbf{分析：}
\begin{itemize}
    \item IDLE 状态：应该只有在 \verb|valid=1| 且无写冲突时才能接收请求
    \item LOOKUP 状态：只有命中且有新的有效请求才能继续接收，缺失时应阻塞
\end{itemize}

\noindent
\textbf{解决：}

\begin{lstlisting}[language=Verilog, caption=正确的 addr_ok 生成]
assign addr_ok = (main_state == MAIN_IDLE && valid && !write_conflict)
               | (main_state == MAIN_LOOKUP && hit && valid && !write_conflict);
\end{lstlisting}


\subsubsection{TLBRD 指令 CSR 冒险未处理}

TLBRD 指令在 WB 级写入 CSR，但 ID 级的 CSR 读指令（CSRRD/CSRXCHG）可能读到旧值。

\noindent
\textbf{现象：}执行 TLBRD 后立即用 CSRRD 读取 TLBEHI，读到的是 TLBRD 执行前的旧值，而非 TLB 表项中的值。

\noindent
\textbf{分析：}之前的实验中未暴露该问题，可能是流水线在相关指令处会停顿，掩盖了冒险现象。使用cache后，流水线变快变紧凑，冒险问题暴露。

已有的冒险检测只处理了 TLBSRCH 写 TLBIDX 的情况：
\begin{lstlisting}[language=Verilog]
wire id_reads_tlbidx = id_is_csr && (id_csr_num == `CSR_TLBIDX);
wire id_stall_tlbidx = id_reads_tlbidx & 
    (ex_pending_tlbsrch | mem_pending_tlbsrch | wb_pending_tlbsrch);
\end{lstlisting}

\noindent
\textbf{解决：}增加 TLBRD CSR 冒险检测，阻塞读取 TLBRD 写入的 CSR：

\begin{lstlisting}[language=Verilog, caption=TLBRD CSR 冒险处理]
// TLBRD writes TLBIDX/TLBEHI/TLBELO0/TLBELO1 in WB. 
// A following CSRRD/CSRXCHG reading these CSRs must wait until TLBRD retires.
wire id_reads_tlbrd_csr = id_is_csr &&
                          ((id_csr_num == `CSR_TLBIDX)  ||
                           (id_csr_num == `CSR_TLBEHI)  ||
                           (id_csr_num == `CSR_TLBELO0) ||
                           (id_csr_num == `CSR_TLBELO1));

wire ex_pending_tlbrd  = u_stage_ex.valid  & u_stage_ex.output_inst_tlbrd;
wire mem_pending_tlbrd = u_stage_mem.valid & u_stage_mem.output_inst_tlbrd;
wire wb_pending_tlbrd  = wb_inst_tlbrd;

wire id_stall_tlbrd_csr = id_reads_tlbrd_csr & 
    (ex_pending_tlbrd | mem_pending_tlbrd | wb_pending_tlbrd);

// 添加到总阻塞条件
wire id_stall = id_stall1 | id_stall2 | id_stall_csr | 
                id_stall_tlbsrch | id_stall_tlbidx | id_stall_tlbrd_csr;
\end{lstlisting}
